\subsection{全SL2 normalizing行列のcurved chart}
\label{subsec:h3-order2-22-curved-SL2}
\begin{theorem}\label{thm:h3-order2-22-curved-SL2}
$F\in\operatorname{SL}(U)$、$U=\langle x,y\rangle$とし、全六関係を
\[
 xy-yx=\alpha u^2+\beta v^2,\quad uv=-vu,\qquad
 ut=F(t)u,\quad vt=F^{-1}(t)v\quad(t=x,y)
\]
とする。
全係数でこれはnoetherian quadratic Koszul AS4整域、$\mu=(+,+,-,-)$である。
この全chartのH3部分は$F=-I,\alpha=\beta=0$の固定符号族だけであり、
全幾何学的矢も既知の符号分類へ帰着する。
\end{theorem}
\begin{proof}
一般$2\times2$行列について$\det F=1$だけを用い、全四三次overlapが消えることを確認する。
正次数を先に比較するmonomial順序で停止し、PBW・Koszul性を得る。
正基底変更で$F$は対角、または$\epsilon J$、$\epsilon=\pm1$、
$J(x)=x,J(y)=y+x$へ移る。
対角型は前節のOre表示を持つ。
$F=\epsilon J$の場合も、負量子平面$T=(uv=-vu)$の上で
\[
 S=T[x;\tau],\quad\tau(u)=\epsilon u,\tau(v)=\epsilon v,
 \qquad B=S[y;\sigma,\delta]
\]
と書ける。ここで
\[
 \sigma(u)=\epsilon u,\ \sigma(v)=\epsilon v,\ \sigma(x)=x,\qquad
 \delta(u)=-xu,\ \delta(v)=xv,\ \delta(x)=-\alpha u^2-\beta v^2
\]
である。関係保存は直接確認でき、二回のOre拡大からnoetherian・整域性・AS4性を得る。

中山がcurvatureを加えても変わらないことを、Koszul topで確認する。
零curvatureの$B(F)$は既にAS4・$\mu=(+,+,-,-)$と分かっている。
そのtopを、正commutatorを$a$、負anticommutatorを$b$として、
偶middleが$a\mathbin{\odot}b$となるよう正規化する。
$h=\alpha u^2+\beta v^2$を加えると、偶middleを
\[
 (a-h)\mathbin{\odot}b
\]
へ変えればよく、奇middleは変わらない。
負主block $-h\mathbin{\odot}b$はcyclicなので、同じ中山のtwisted cyclic tensorが得られる。
偶階数は2、奇階数は4のままで、一対三flatteningも非退化である。
PBW・Koszulから双対Hilbert値は$14641$であり、このtensorが完全なFrobenius対合を与える。
従って中山は元の$(+,+,-,-)$に等しい。
ここでは符号自己同型が逆写像と同じなので、右左規約による反転も結論を変えない。

対角型の全H3は前節までに認識済みである。
非対角二型を考え、正三次空間を
\[
 E=\langle x^3,x^2y,xy^2,y^3,xu^2,yu^2,xv^2,yv^2\rangle,
 \qquad E'=\langle xuv,yuv\rangle
\]
に分ける。
正shear $J$は負変数を固定して$B$の自己同型となる。
$E$はその作用で不変で、正次数filtrationの随伴次数上では
$\operatorname{Sym}^3J\oplus J\oplus J$となる。
従って固定核は高々3次元である。
実際に$x^3,xu^2,xv^2$は固定されるので
\[
 E^J=\langle x^3,xu^2,xv^2\rangle
\]
である。この等式は全curvatureで、$J-I$の定数minor 6によっても検証できる。
また$E'$の固定核は$\langle xuv\rangle$である。

$u$の正常化作用は$E$上で$\epsilon J$、$E'$上で$-\epsilon J$である。
$Gu=-uG$が必要なので、$\epsilon=1$では$G\in\langle xuv\rangle$、
$\epsilon=-1$では$G\in\langle x^3,xu^2,xv^2\rangle$に限られる。
いずれも$x$を一次因子に持ち、整域商を与えない。
これで全非対角型も排除し、既知符号族だけが残る。
\end{proof}

検証器は
\[
 \texttt{verify\_cubic\_order2\_22\_curved\_SL2.py}
\]
である。$\det F=1$の多項式イデアルを法とした全PBW overlap、
cyclic負topと偶階数の定数minor、両Jordan型でのshear核の定数minor 6を再生する。
一般の非逆行列対や、退化した正二次射影を覆うとは主張していない。
